% ============================================================================
% CHAPTER 2 - THEORETICAL FOUNDATION
% File: docs/thesis/theory/ch2_renewable_potential_foundations.tex
% Sections: (1) potential definitions and capacity accounting,
%           (2) wind resource and turbine performance,
%           (3) PV thermal performance,
%           (4) levelized cost and cost-basis consistency.
% Operational values live in docs/thesis/methods/group1_technical_potential.tex
% and docs/thesis/methods/group2_core_costs.tex
% ============================================================================

% ----------------------------------------------------------------------------
\section{Spatially Explicit Renewable-Energy Potential}
\label{sec:renewable_potential_foundations}
% ----------------------------------------------------------------------------

\subsection{Definitions of Potential}
\label{subsec:potential_definitions}

Renewable-energy assessments distinguish the physical resource, the
geographical potential, the technical potential, the economic potential
and feasible deployment. The physical resource is the energy carried by
the natural flux. The geographical potential is the share of that
resource located on land that is physically and legally available.
The technical potential converts the geographical potential into
electricity using a specified technology, its performance and its
land requirements. The economic potential is the share of the technical
potential that satisfies an economic criterion, such as a maximum
levelized cost of electricity (LCOE). Feasible deployment adds
institutional, social and market constraints that are rarely mappable
\cite{heinrichs2026,mckenna2022,turkovska2024}.

These categories are not interchangeable, and published estimates for
the same country can differ by large factors because studies combine
different exclusion criteria, technology assumptions, land-use factors
and economic thresholds. Reviews of solar and onshore wind assessments
attribute much of this spread to definitions and reporting rather than
to the resource itself \cite{mckenna2022,turkovska2024}. Any comparison
of a modeled potential with a published figure is therefore meaningful
only when both use the same definition and unit.

\subsection{Capacity and Energy Accounting}
\label{subsec:capacity_accounting}

Let $i$ index a grid cell and $j$ a technology. The eligible area
$A^{\mathrm{elig}}_{i,j}$ (km$^2$) is the area that remains after the
mapped exclusions. Installable capacity follows from two factors: a
power density $\mathrm{PD}_j$ (MW/km$^2$), which states how much
capacity fits on one square kilometre of project area, and a
land-utilization factor $\mathrm{LUF}_j\in[0,1]$, which states the share
of the eligible area that is expected to host projects:
\begin{equation}
    C_{i,j} = A^{\mathrm{elig}}_{i,j}\,\mathrm{LUF}_j\,\mathrm{PD}_j .
    \label{eq:theory_capacity}
\end{equation}
Annual energy follows from the capacity factor:
\begin{equation}
    E_{i,j} = 8760\,C_{i,j}\,\mathrm{CF}_{i,j}
    \quad [\mathrm{MWh/yr}] .
    \label{eq:theory_energy}
\end{equation}

Two properties of Equation~\ref{eq:theory_capacity} matter for the
interpretation of results. First, only the product
$\mathrm{LUF}_j\,\mathrm{PD}_j$ affects capacity, so the two factors
cannot be identified separately from capacity outcomes alone. Second,
because both factors are uniform within a country and technology, they
scale the potential of every cell by the same amount and leave cost per
unit of energy, and therefore the ranking of cells by LCOE, unchanged.
They govern quantities expressed in GW or TWh/yr, not the relative
merit of locations.

Capacity must be expressed on one electrical basis throughout. For PV,
land requirements are set by the direct-current (DC) array, so power
density is naturally expressed in MW$_{\mathrm{dc}}$/km$^2$
\cite{bolinger2022}. The cost per unit of capacity must then refer to the
same basis; the inverter loading ratio
$\mathrm{ILR}=C_{\mathrm{dc}}/C_{\mathrm{ac}}$ links the two bases when a
source reports alternating-current (AC) ratings.

\subsection{The Residual Land-Utilization Factor}
\label{subsec:luf_theory}

In the literature, land-use, suitability or availability factors often
bundle together several reasons why eligible land is not built on:
protected areas, slope, land cover, distance to settlements, land tenure,
competing uses and social acceptance. When some of these constraints are
already removed by explicit spatial exclusions, applying a bundled
factor afterwards removes the same land twice and biases the potential
downwards.

GeoFREA therefore defines the LUF as a residual factor. It represents
only constraints that the exclusion layers do not encode: land-tenure
fragmentation, competing uses not mapped as exclusions, micro-siting
losses inside otherwise eligible land, permitting and social acceptance.
A literature value is admissible as a LUF only if its definition allows
the mapped constraints to be separated out. Because few sources satisfy
this condition, residual LUF values are typically judgment-based and
must be reported as such.

\subsection{Power Density and Footprint Definitions}
\label{subsec:power_density_theory}

Power density depends on the denominator. For utility-scale PV, the
direct array area includes module rows and inverter pads, whereas the
fenced or leased site also includes internal roads, setbacks and
undevelopable patches. In a large empirical sample of United States
plants built through 2019, Bolinger and Bolinger measured density on the
direct array area and expressed it in DC terms, found fixed-tilt plants
to be denser than single-axis trackers because of their higher ground
coverage ratio, and showed that density rose markedly over the decade as
module efficiency improved \cite{bolinger2022}. Earlier benchmarks were
measured on total project area, in AC terms and on an older fleet
\cite{ong2013}, and national technical-potential studies adopt a single
assumed density \cite{lopez2012}. Densities from older studies or from a
different footprint definition are therefore not transferable without
conversion \cite{turkovska2024}.

For onshore wind, the relevant denominator is the spacing envelope of
the turbine array, not the much smaller area directly disturbed by
foundations, roads and substations
\cite{harrisonatlas2022,mckenna2022}. Most of the envelope remains
available for agriculture or grazing. With turbines spaced $s_x D$
across and $s_y D$ along the prevailing wind, where $D$ is the rotor
diameter, the geometric density is
\begin{equation}
    \mathrm{PD}_{\mathrm{wind}}
    = \frac{P_{\mathrm{rated}}}{s_x s_y D^2}
    = \frac{\pi}{4\,s_x s_y}\,\mathrm{SP},
    \label{eq:theory_wind_pd}
\end{equation}
where $\mathrm{SP}=P_{\mathrm{rated}}/(\pi D^2/4)$ is the specific power
of the turbine. In consistent units, SP in W/m$^2$ equals MW/km$^2$, so
the geometric density depends only on the specific power and on the
spacing product $s_x s_y$. Low-specific-power turbines, typical of
low-wind sites, yield lower capacity per unit area for the same spacing.
Empirical densities of real wind farms vary widely with terrain, layout
and the footprint definition adopted: convex hulls, Voronoi cells and
buffers around turbines give different areas for the same plant
\cite{harrisonatlas2022,turkovska2021,turkovska2024}. Geometric and
empirical densities are therefore alternative parameterizations, useful
as mutual checks, not repeated observations of one constant.

\subsection{Technology Vintage and Climate Period}
\label{subsec:technology_vintage}

Land requirements, performance and costs change with technology
generation. Combining a power density of one generation with costs of
another produces an inconsistent plant. A technology reference year
$t_{\mathrm{ref}}$ is therefore used to fix the vintage of every
technology parameter (module type, turbine, hub height, densities and
costs).

The technology vintage and the climate period are distinct. A
climate-robustness assessment can hold the technology fixed and expose
it to the reference and projected climates, so that differences across
climate members reflect climate alone and not undocumented changes in
equipment. The price of this design is that the evaluated plant is a
fixed-vintage counterfactual: its technology does not evolve with the
projected climate window, and a plant commissioned at
$t_{\mathrm{ref}}$ would only partly operate within a mid-century
window. Both points must be stated when results are interpreted.

\subsection{Evidence Quality and Transferability}
\label{subsec:evidence_quality}

A framework meant to be transferable across countries must separate the
equations, which are general, from the parameter values, which are
calibrated per country and technology. Transferability of the equations
does not imply transferability of the values. Whether a value measured
in one market can represent another depends on whether the technology,
the footprint or cost definition, the vintage and the units match.
Parameter provenance must therefore be recorded together with the value,
and values derived from judgment must remain distinguishable from those
measured locally \cite{heinrichs2026,mckenna2022}.

Parameter estimation must also stay independent of validation. Plant
inventories used to derive a density or a cost should not be reused as
independent evidence that the same workflow places plants correctly.

% ----------------------------------------------------------------------------
\section{Wind Resource and Turbine Performance}
\label{sec:wind_performance_theory}
% ----------------------------------------------------------------------------

\subsection{Weibull Representation of the Wind Resource}
\label{subsec:wind_weibull}

When long hourly series are not used, the wind-speed distribution at a
site is commonly represented by a two-parameter Weibull density
\begin{equation}
    f(u;A,k_{\mathrm{W}})
    = \frac{k_{\mathrm{W}}}{A}\left(\frac{u}{A}\right)^{k_{\mathrm{W}}-1}
      \exp\!\left[-\left(\frac{u}{A}\right)^{k_{\mathrm{W}}}\right],
    \qquad u\ge 0,
    \label{eq:wind_weibull_pdf}
\end{equation}
with scale $A$ (m/s) and dimensionless shape $k_{\mathrm{W}}$. The mean
wind speed is
\begin{equation}
    \bar{u} = A\,\Gamma\!\left(1+\frac{1}{k_{\mathrm{W}}}\right).
    \label{eq:wind_weibull_mean}
\end{equation}
The Weibull form describes the marginal distribution only. It carries
no information on temporal sequence, seasonality or co-occurrence with
other variables, which limits its use in chronological system models
but not in annual yield estimation. The Global Wind Atlas version 3
provides mean wind speed and Weibull parameters at 10, 50, 100, 150 and
200~m for the period 2008--2017, obtained by downscaling ERA5 with a
mesoscale model at about 3~km and a microscale flow model to about 250~m
\cite{davis2023globalwindatlas,dtuwind2019gwa3release,gwaV3dataset}. It
also provides an air-density layer derived from a different reanalysis
(CFSR) interpolated to the atlas elevation \cite{gwaV3dataset}. Its
capacity-factor layers are gross values computed at 100~m for three
commercial turbines representing IEC classes I to III, without wake or
other losses \cite{davis2023globalwindatlas}.

\subsection{Hub Height}
\label{subsec:wind_hub_height_theory}

Wind speed increases with height in a way that depends on surface
roughness, stability and terrain. Extrapolation with a power law
$u(h_2)=u(h_1)(h_2/h_1)^{\alpha}$ requires a shear exponent $\alpha$
that is neither constant in space nor well constrained over large
domains. Taking the resource directly from a dataset layer at or near
the hub height avoids this extra parameter, at the cost of a small,
documented mismatch when the turbine hub does not coincide with an
available layer.

\subsection{Design Classes and Power Curves}
\label{subsec:wind_classes_curves}

IEC 61400-1 assigns turbines to design classes defined by a 50-year
reference wind speed $V_{\mathrm{ref}}$ and an associated annual mean
wind speed $V_{\mathrm{ave}}$ at hub height, together with turbulence
categories \cite{iec6140012025,larsen2022}. Table~\ref{tab:wind_iec_classes}
lists the reference values. The current edition also defines a class
for tropical-cyclone regions and site-specific class S designs.

\begin{table}[htbp]
    \centering
    \caption{Reference wind speeds of the IEC 61400-1 design classes.}
    \label{tab:wind_iec_classes}
    \small
    \begin{tabular}{lcc}
        \toprule
        Class & $V_{\mathrm{ref}}$ (m/s) & $V_{\mathrm{ave}}$ (m/s) \\
        \midrule
        I   & 50.0 & 10.0 \\
        II  & 42.5 & 8.5  \\
        III & 37.5 & 7.5  \\
        S   & site-specific & site-specific \\
        \bottomrule
    \end{tabular}
\end{table}
% verify: values of the tropical-cyclone class (T) in the edition cited; not tabulated here until checked

Design classes are structural: they certify that a turbine survives the
loads of a wind climate. They do not, by themselves, determine energy
yield. Yield is set by the power curve, which in turn is shaped mainly by
specific power. The market pairs low-wind sites with large rotors and
low specific power, so annual mean speed is only a screening indicator
of the turbine that would actually be installed. Formal class assignment
also requires extreme-wind and turbulence information that spatial
atlases do not provide.

For reproducibility, power curves should come from openly documented
reference turbines whose provenance and licence are known, such as the
IEA Wind Task 37 land-based reference turbine \cite{bortolotti2019}.
Software libraries that distribute manufacturer curves do not, by being
open source, make those curves open data.

\subsection{Gross Capacity Factor}
\label{subsec:wind_gross_cf}

For a power curve $P(u)$ evaluated at the site air density, the gross
capacity factor over a Weibull distribution is
\begin{equation}
    \mathrm{CF}_{\mathrm{gross}}
    = \frac{1}{P_{\mathrm{rated}}}
      \int_{u_{\mathrm{in}}}^{u_{\mathrm{out}}} P(u)\,f(u;A,k_{\mathrm{W}})\,du ,
    \label{eq:wind_gross_cf}
\end{equation}
where $u_{\mathrm{in}}$ and $u_{\mathrm{out}}$ are the cut-in and
cut-out speeds.

\subsection{Air Density}
\label{subsec:wind_air_density}

The kinetic power flux through a rotor is proportional to
$\tfrac{1}{2}\rho u^3$. Reference power curves are stated at a reference
density $\rho_0$, usually $1.225$~kg/m$^3$. Warm and elevated sites have
lower density and, below rated speed, produce less power at the same
wind speed. For dry air,
\begin{equation}
    \rho = \frac{p}{R_d\,T},
    \label{eq:wind_density_dry}
\end{equation}
with $p$ in Pa, $T$ in K and $R_d=287.05$~J\,kg$^{-1}$\,K$^{-1}$. Humidity
lowers density slightly further; the effect is small relative to that of
temperature and elevation.

For pitch-regulated turbines, IEC 61400-12-1 normalizes power curves to a
reference density by rescaling wind speed \cite{iec61400122022}. The
inverse operation gives the power at site density from a reference curve
$P_0$:
\begin{equation}
    P(u,\rho) = P_0\!\left(u\left(\frac{\rho}{\rho_0}\right)^{1/3}\right).
    \label{eq:wind_density_correction}
\end{equation}
This rescaling already accounts for density; multiplying the output by
$\rho/\rho_0$ as well would count the effect twice.

\subsection{Gross-to-Net Losses}
\label{subsec:wind_losses_theory}

Net output is lower than gross output because of wake interaction inside
the array, turbine unavailability, electrical losses in the collection
system and degradation of blade performance by soiling, erosion or
icing \cite{leefields2021}. When component losses $\ell_m$ act as
successive reductions, the retained fraction of gross energy is
\begin{equation}
    \eta = \prod_{m}(1-\ell_m),
    \qquad
    \mathrm{CF}_{\mathrm{net}} = \eta\,\mathrm{CF}_{\mathrm{gross}},
    \label{eq:wind_net_cf}
\end{equation}
and the combined loss is $1-\eta$. Effects already represented in the
power curve or in the density correction must not be included again in
$\eta$.
Curtailment is a property of the power system, not of the plant, and
long-term degradation is a separate time-dependent process.

% ----------------------------------------------------------------------------
\section{Thermal Performance of Photovoltaic Modules}
\label{sec:pv_thermal_theory}
% ----------------------------------------------------------------------------

\subsection{Temperature Dependence of Module Power}
\label{subsec:pv_temperature}

PV nameplate power is rated at standard test conditions: 1000~W/m$^2$,
AM1.5 spectrum and 25\,$^\circ$C cell temperature. In operation, cells
run well above ambient temperature under strong irradiance, and power
falls almost linearly with temperature:
\begin{equation}
    f_T = 1+\gamma\,(T_{\mathrm{cell}}-25\,^\circ\mathrm{C}),
    \label{eq:pv_thermal_factor}
\end{equation}
where $\gamma$ is the temperature coefficient of maximum power, in
K$^{-1}$; datasheets report it in \%/K, so $-0.30$\,\%/K corresponds to
$\gamma=-0.0030$~K$^{-1}$. It must not be confused with the coefficients
for open-circuit voltage or short-circuit current. For crystalline
silicon, $\gamma$ is negative and its magnitude decreases from p-type
PERC to n-type TOPCon and heterojunction cells
\cite{bamisile2025,itrpv2025}.

Cell temperature models express $T_{\mathrm{cell}}$ as ambient
temperature plus a heating term that grows with irradiance and falls with
wind speed, as in the model of Faiman,
$T_{\mathrm{mod}} = T_{\mathrm{amb}} + G_{\mathrm{POA}}/(U_0+U_1 w)$
\cite{faiman2008}. At equal irradiance and wind, a change in ambient
temperature therefore shifts cell temperature by the same amount.

\subsection{Climate Perturbation of a Baseline Yield}
\label{subsec:pv_climate_perturbation}

Gridded specific-yield datasets give the long-term energy output per
unit of DC capacity of a reference PV configuration, already net of
thermal and system losses under the reference climate. A projected
climate can be represented as a perturbation of that baseline. With an
irradiance ratio $\delta_G$ and an ambient warming $\Delta T$,
\begin{equation}
    Y_c = Y_{\mathrm{ref}}\,\delta_G\,
      \frac{1+\gamma\,(\bar{T}_{\mathrm{cell,ref}}+\Delta T-25)}
           {1+\gamma\,(\bar{T}_{\mathrm{cell,ref}}-25)}
    \approx Y_{\mathrm{ref}}\,\delta_G\,(1+\gamma\,\Delta T),
    \label{eq:pv_perturbation}
\end{equation}
where the approximation drops terms of order $\gamma^2$ and equates the
change in cell temperature with the change in ambient temperature. It
assumes that warming is spread evenly over productive hours, ignores the
effect of irradiance and wind changes on cell temperature and uses annual
means. Its error is small relative to the irradiance term, but it does
not capture changes in the diurnal or seasonal timing of warming.

Technology evolution, such as a smaller $|\gamma|$, also changes the
thermal response but is not a climate effect. Keeping the module type
fixed across climate members attributes differences between members to
climate only.

% ----------------------------------------------------------------------------
\section{Levelized Cost and Cost-Basis Consistency}
\label{sec:lcoe_theory}
% ----------------------------------------------------------------------------

\subsection{Levelized Cost of Electricity}
\label{subsec:lcoe_definition}

For a plant with overnight capital cost $I$ per unit of capacity, fixed
operating cost $F$ per unit of capacity and year, real discount rate $r$,
lifetime $n$ and annual energy $E$ per unit of capacity, the levelized
cost under constant real annual costs is
\begin{equation}
    \mathrm{LCOE} = \frac{I\,\mathrm{CRF}(r,n) + F}{E},
    \qquad
    \mathrm{CRF}(r,n) = \frac{r(1+r)^n}{(1+r)^n-1}.
    \label{eq:theory_lcoe}
\end{equation}
The expression is consistent only if costs, discount rate and energy
share one convention: costs in constant currency with a real rate, or
nominal costs with a nominal rate. A nominal rate $i$ is converted to a
real rate with the Fisher relation
\begin{equation}
    1+r = \frac{1+i}{1+\pi},
    \label{eq:theory_fisher}
\end{equation}
where $\pi$ is the expected inflation over the asset life. Because
renewable plants are capital-intensive, the LCOE is more sensitive to
$r$ than to most technical parameters, and international differences in
financing costs explain a large share of observed differences in LCOE
\cite{irena2023financing}.

\subsection{Cost Basis and Scope}
\label{subsec:cost_basis}

A unit cost is meaningful only with its denominator and scope. For PV,
cost per kW$_{\mathrm{ac}}$ and per kW$_{\mathrm{dc}}$ differ by the
inverter loading ratio: for the same total plant cost,
\begin{equation}
    c_{\mathrm{dc}} = \frac{c_{\mathrm{ac}}}{\mathrm{ILR}} .
    \label{eq:theory_capex_basis}
\end{equation}
The ILR depends on plant design and should be documented rather than
assumed universal \cite{nrelATB2024PV}. The scope matters equally:
reported total installed costs may include grid connection, land and
development costs, and operating costs may be reported as full-service
contract prices or as all-in expenditure. Combining a cost that already
includes grid connection with an explicit connection cost counts that
cost twice.

Capacity factors also have more than one definition. The capacity factor
of plants commissioned in a given year is not the same quantity as the
fleet capacity factor $E_{\mathrm{annual}}/(8760\,\bar{P})$ computed
from national generation and average installed capacity, which mixes
vintages and includes curtailment and outages. Neither is a long-term
climatological capacity factor of a candidate site; both are useful only
as plausibility references.
